\subsection{The quasi-\texorpdfstring{$F$}{F}-splitting matrix and complexity}
\label{subsec:quasi-height-complexity}

The full-height iteration used here is an existing algorithm.
Kawakami, Takamatsu, and Yoshikawa give its exact criterion
\citep[Theorem~C, version~3, p.~4; Theorem~5.8, pp.~37--39]{KTY2025};
the criterion was already present in version~2 of their preprint, p.~5.
Batubara, Garzella, and Pan identify the fixed finite-dimensional operator
and implement its matrix iteration
\citep[Section~3.4 and Algorithm~4, p.~7]{BGP2025}.
Yobuko identifies quasi-$F$-split height with Artin--Mazur height for
Calabi--Yau varieties
\citep[Theorem~4.5; author version, pp.~9--10]{Yobuko2019}.
We analyze this algorithm under the circuit input convention of this paper.

Put $S=\mathbb F_p[x_1,\ldots,x_d]$ and write
\[
 D=d(p-1),\qquad g=F^{p-1},\qquad
 V=S_D,\qquad a=(p-1,\ldots,p-1),
\]
\[
 W=\dim_{\mathbb F_p}V=\binom{dp-1}{d-1}.
\]
Thus $V=S_D$ is the space of homogeneous degree-$D$ polynomials.
Choose a homogeneous integral coefficient lift
$G=\sum_{|\beta|=D}\widetilde c_\beta x^\beta$ of $g$, and define
\begin{equation}\label{eq:quasi-delta}
 \Delta=\left(
 \frac{G^p-\sum_{|\beta|=D}\widetilde c_\beta^{\,p}x^{p\beta}}p
 \right)\bmod p.
\end{equation}
This is the monomial representative $\Delta_1(g)$ of the cited criterion;
see also \citet[Proposition~3.2, Remark~3.4 and Algorithm~1, p.~5]{BGP2025}.
The sum is over distinct collected monomials.
The numerator is divisible by $p$ coefficientwise. Moreover, replacing
$G$ by $G+pH$, or a scalar lift by $\widetilde c_\beta+pb$, changes the
corresponding $p$th power by a multiple of $p^2$. Thus
\eqref{eq:quasi-delta} is independent of the chosen integer coefficient
lifts. This includes $p=2$ and requires no Teichm\"uller representatives.
Computationally, form the numerator modulo $p^2$, choose each residue in
$[0,p^2)$, and divide it by $p$; the resulting residue in $[0,p)$ is exact.

Define the $\mathbb F_p$-linear map $U$ on monomials by
\[
 U(x^\gamma)=
 \begin{cases}
 x^{(\gamma-a)/p},&\gamma\equiv a\pmod p,\\
 0,&\text{otherwise}.
 \end{cases}
\]
Since $\deg\Delta=pD$, the map
\[
 M:V\longrightarrow V,\qquad v\longmapsto U(\Delta v)
\]
preserves $V$: its retained monomials have degree
$(pD+D-d(p-1))/p=D$. In the monomial basis,
\begin{equation}\label{eq:quasi-matrix-entry}
 M_{\alpha,\beta}=[x^{p\alpha+a-\beta}]\Delta,
 \qquad |\alpha|=|\beta|=D,
\end{equation}
where a negative exponent gives zero. Let $\ell(v)=[x^a]v$.
The only degree-$D$ monomial outside $(x_1^p,\ldots,x_d^p)$ is $x^a$.
Consequently the cited height criterion gives
\begin{equation}\label{eq:quasi-height-iteration}
 h(X)=\inf\{j\ge1:\ell(M^{j-1}g)\ne0\},\qquad \inf\varnothing=\infty.
\end{equation}

For the field hypotheses, apply Kawakami, Takamatsu, and Yoshikawa's
projective criterion and Yobuko's
comparison after extension to $\overline{\mathbb F}_p$. The Cartier map
over a perfect field takes $p$th roots of coefficients, but every
coefficient of this recurrence lies in $\mathbb F_p$, so the recurrence
is exactly the linear one displayed above. Formal-group height is
unchanged by extending the base field, since nonvanishing of the first
nonzero coefficient of the $p$-series is unchanged. The adjunction and
cohomology calculations already made verify the Calabi--Yau hypotheses.
This applies also when $d=3$.

\begin{lemma}[A matrix stopping bound]
\label{lem:quasi-matrix-stopping}
In \eqref{eq:quasi-height-iteration}, a finite height is at most $W$.
Thus testing $j=1,\ldots,W$ detects infinite height as well. For the
smooth inputs here one may use the smaller cutoff $\min(W,B_d)$.
\end{lemma}
\begin{proof}
The sequence $y_r=\ell(M^rg)$ satisfies a linear recurrence of order
at most $W$ by Cayley--Hamilton. If its first $W$ terms vanish, the
recurrence forces every subsequent term to vanish. Otherwise the first
nonzero term has index at most $W-1$, which proves the assertion by
\eqref{eq:quasi-height-iteration}. This argument requires no computation
of the characteristic polynomial. The alternative bound $B_d$ is
\eqref{eq:HC1}.
\end{proof}

Write $\mathrm{FPSPACE}$ for function problems computable deterministically
with a number of work bits polynomial in the input length.

\begin{theorem}[Resource bounds for exact height]
\label{thm:quasi-height-resource-bounds}
For a binary-encoded prime $p$ and a division-free circuit of size $s$
in $d\ge3$ explicitly listed variables, promised to define a smooth
degree-$d$ hypersurface in $\mathbb P^{d-1}$, exact Artin--Mazur height,
including $\infty$, belongs to $\mathrm{FPSPACE}$.
There is also a deterministic dense implementation with bit complexity
\[
 2^{O(d\log p)}\operatorname{poly}(s,d,\log p).
\]
The polynomial-space and dense-time bounds refer to separate
implementations of the same height criterion.
\end{theorem}
\begin{proof}
Let $L$ denote input length and set $E=pD=dp(p-1)$. All scalar
calculations below take place modulo $p$ or $p^2$, so each scalar has
$O(\log p)$ bits. A monomial of total degree at most $E$ has an exponent
vector using $O(d\log(E+1))$ bits. Furthermore,
\[
 \log W=O(d(1+\log p)),\qquad
 \log(E+1)=O(\log d+\log p).
\]
These lengths are polynomial in $L$, because $d$ is explicitly bounded
by input length.

\paragraph{Coefficient access.}
For a degree-$d$ output coefficient of $F$, recurse through its arithmetic
circuit. At a product gate, enumerate all coordinatewise splits of the
target exponent vector, obtain the two child coefficients, and accumulate
their product modulo $p$. At a sum gate, add the child coefficients.
Constants and variables are immediate. A recursion path has at most $s$
gates, and each frame stores polynomial-length exponent vectors, gate
addresses, a loop counter and an $O(\log p)$-bit accumulator. Intermediate
gates need not be homogeneous. Equivalently, this computation evaluates
the circuit modulo all monomials of total degree greater than $d$, which
retains the required output exactly. The degree-$d$ output is selected
using the promise; an output coefficient of any other total degree is zero.

Coefficients of a power $A^r$ are obtained by the same convolution
recursion applied to
\[
 A^r=A^{\lfloor r/2\rfloor}A^{\lceil r/2\rceil},
\]
with base cases $r=0,1$. The power recursion has depth
$O(\log(r+1))$. First compute coefficients of $g=F^{p-1}$ modulo $p$.
Treat their representatives in $[0,p)$ as the coefficients of a virtual
homogeneous lift $G$, and compute coefficients of $G^p$ modulo $p^2$.
The two powering recursions add only $O(\log p)$ depth to the original
circuit recursion. In the pure-power sum in \eqref{eq:quasi-delta}, a
requested exponent has at most one contributing term, determined by
coordinatewise divisibility by $p$. Scalar powering modulo $p^2$ and
the exact scalar division by $p$ complete the coefficient calculation.
Thus every coefficient of $\Delta$, and hence every matrix entry in
\eqref{eq:quasi-matrix-entry}, is accessible in polynomial space.

\paragraph{Matrix powers.}
For $r\ge2$, compute one entry recursively from
\[
 (M^r)_{\alpha,\beta}
   =\sum_{|\gamma|=D}
      (M^{\lfloor r/2\rfloor})_{\alpha,\gamma}
      (M^{\lceil r/2\rceil})_{\gamma,\beta}\pmod p.
\]
The cases $r=0,1$ use the identity and the entry oracle just constructed.
For $r<W$, the matrix recursion has depth $O(\log W)$.
Every frame stores two index vectors, the current summation index, a
binary exponent, and one field accumulator. All have polynomial length.
An entry-oracle call adds only the polynomial stack used for coefficients.
Finally, compute
\[
 \ell(M^{j-1}g)=
 \sum_{|\beta|=D}(M^{j-1})_{a,\beta}g_\beta\pmod p
\]
and scan $j=1,\ldots,W$ using a binary counter. Return the first nonzero
test, or $\infty$ if all vanish. Lemma~\ref{lem:quasi-matrix-stopping}
proves correctness. The outer scan does not create a recursion chain of
length $W$; every requested power uses divide and conquer. The complete
stack, counters, exponents and scalar accumulators use $L^{O(1)}$ space.

\paragraph{A separate dense time bound.}
The number of monomials of total degree at most $E$ is
\[
 V_E=\binom{E+d}{d}
 \le \bigl(e[p(p-1)+1]\bigr)^d
 \le (ep^2)^d,
 \qquad W\le\binom{dp}{d}\le(ep)^d.
\]
Expand the original circuit modulo $p$ through total degree $d$,
using at most $\binom{2d}{d}\le4^d$ positions. Dense binary powering
computes $g$ and $G^p$ in $O(\log p)$ polynomial-product stages.
Naive coefficient convolution uses at most $V_E^2$ scalar products
per stage. Form $\Delta$ by \eqref{eq:quasi-delta} and construct the
$W^2$ matrix entries by coefficient lookup. At most $W-1$ matrix-vector
products then test the whole sequence, using $O(W^3)$ field operations.
Degree indexing and scalar arithmetic contribute a fixed polynomial in
$d,\log p$; even sequential array-access implementations contribute
only a fixed polynomial in $V_E,W$. Thus the total bit complexity is
a fixed polynomial in $s,d,\log p,V_E,W$, giving
$2^{O(d\log p)}\operatorname{poly}(s,d,\log p)$ uniformly in all
parameters. This dense implementation uses exponentially large arrays,
whereas the preceding recursion recomputes entries to obtain the
polynomial-space bound.
\end{proof}

